\begin{proposition}[The theta surface and its resolution]
\label[proposition]{prop:theta-resolution}
The theta section $X\subset P$ is an integral normal symmetric ample Cartier
divisor with $\Sing(X)=\{0\}$.  The morphism $\rho:S\to X$ of
\cref{lem:norm-fibre} is its minimal resolution and is an isomorphism over
$X\setminus\{0\}$.  Its exceptional locus is the smooth elliptic curve
$E_{\exc}$ with $E_{\exc}^2=-2$; in particular, $(X,0)$ is a simple elliptic
singularity of type $\wt E_7$.
\end{proposition}

\begin{proof}
A basis of $H^0(C,K_C)$ is represented on the affine chart by
\[
 \frac{dx}{y},\qquad
 \frac{dx}{wy},\qquad
 \frac{dx}{w^2y},\qquad
 \frac{x\,dx}{w^2y}.
\]
Indeed these forms are regular both at the ramification points and at the
three points above $O$.  After multiplication by the common factor
$w^2y/dx$, the canonical map is
\[
 (x,w)\longmapsto[w^2:w:1:x].
\]
The map is birational, so $C$ is not hyperelliptic.  Among the ten
quadratic monomials in the four canonical coordinates there is the evident
relation $X_0X_2-X_1^2=0$, and it is the only one.  After restriction to
$C$, the other nine monomials may be represented as
\[
 \underbrace{1,x,Y,x^2}_{\tau\text{-eigenvalue }1};\qquad
 \underbrace{w,xw,wY}_{\tau\text{-eigenvalue }\zeta_3};\qquad
 \underbrace{w^2,xw^2}_{\tau\text{-eigenvalue }\zeta_3^2}.
\]
At any of the three points above $O$ their pole orders within these three
eigenspaces are, respectively,
\[
 (0,2,3,4),\qquad (1,3,4),\qquad (2,4).
\]
They are therefore linearly independent within each eigenspace, and terms
from distinct eigenspaces cannot cancel.  Thus the canonical ideal has a
single quadric,
\[
 X_0X_2-X_1^2=0,
\]
which has rank three and one ruling.  By geometric Riemann--Roch, every
$g^1_3$ on a nonhyperelliptic genus-four curve determines a ruling of the
canonical quadric.  Hence $\kappa$ is the unique $g^1_3$ on $C$.
Riemann--Kempf \cite{Kempf1973} shows that the Abel map is an isomorphism
over the locus $h^0=1$.  The corresponding points of $N$ are smooth by the
critical-point calculation above, and $|\kappa|$ is the only
positive-dimensional Abel fibre.  It follows that
$\Sing(X)\subset\{0\}$.

The Riemann singularity theorem \cite{Kempf1973} gives
$\operatorname{mult}_0(\Theta_\kappa)=h^0(C,\kappa)=2$.  The restricted
theta equation is nonzero, since $\dim X=2<\dim P$, and lies in
$\mathfrak m_{P,0}^2$.  Thus $X$ is singular at the origin, and
\begin{equation}
 \Sing(X)=\{0\}.
\label{eq:singular-locus}
\end{equation}
The Cartier divisor $X$ is generically reduced and Cohen--Macaulay, hence
reduced.  Since $L$ is ample, $X$ is connected.  Suppose that $X$ were
reducible.  Two irreducible components joined by an edge in the incidence
graph must then meet.  Both have dimension two in the smooth threefold $P$,
so every nonempty intersection of two distinct components has dimension at
least one.  At such a point the local ring of $X$ has at least two minimal
primes and is therefore not regular.  Hence the intersection is contained in
$\Sing(X)=\{0\}$, contradicting its positive dimension.  Thus $X$ is
integral.  Being Cohen--Macaulay, it satisfies $S_2$, while
\eqref{eq:singular-locus} gives regularity in codimension one.  Serre's
criterion makes $X$ normal.

The uniqueness of the $g^1_3$ shows that the Abel map is an isomorphism away
from $\Gamma$.  Hence $\rho$ is a resolution whose only exceptional curve is
$E_{\exc}$.  Since this curve is elliptic with self-intersection $-2$, the
resolution is minimal and the simple-elliptic classification identifies the
germ with type $\wt E_7$ \cite{Laufer1977,Pinkham1977}.
\end{proof}

\begin{lemma}[Surface invariants]\label[lemma]{lem:surface-invariants}
The minimal resolution $S$ satisfies
\[
 K_S^2=16,\qquad \chi(\mathcal O_S)=2,\qquad e(S)=8.
\]
\end{lemma}

\begin{proof}
From $K_N\equiv a^*L$ and $X\in|L|$ one obtains
\[
 K_N^2=(a^*L)^2=L^2\cdot[X]=L^3=18.
\]
For the blowup of the $\wt E_6$ vertex,
$K_{\widehat N}=\beta^*K_N-E_0$, so
\[
 K_{\widehat N}^2=18+E_0^2=15.
\]
The contraction $\widehat N\to S$ is the inverse of a blowup at a smooth
point, and therefore $K_S^2=16$.

The Cartier sequence
\[
 0\longrightarrow L^{-1}\longrightarrow\mathcal O_P
 \longrightarrow\mathcal O_X\longrightarrow0
\]
gives $\chi(\mathcal O_X)=3$: on the abelian threefold $P$,
$\chi(\mathcal O_P)=0$, $\chi(L)=3$, and
$\chi(L^{-1})=-\chi(L)=-3$.  A simple elliptic singularity has local
geometric genus one; equivalently,
$R^1\rho_*\mathcal O_S$ is a skyscraper sheaf of length one
\cite{Laufer1977}.  Since $X$ is normal,
$\rho_*\mathcal O_S=\mathcal O_X$, and Leray gives
\[
 \chi(\mathcal O_S)=\chi(\mathcal O_X)-1=2.
\]
Noether's formula now yields
\[
 e(S)=12\chi(\mathcal O_S)-K_S^2=24-16=8.
\]
\end{proof}

\begin{corollary}[The intersection complex]
\label[corollary]{cor:theta-ic}
There is a decomposition
\[
 R\rho_*\QQ_S[2]\simeq\IC_X\oplus\delta_0.
\]
Moreover,
\[
 \chi(\IC_X)=7,
 \qquad
 \Stab_P(X)=\{0\}.
\]
\end{corollary}

\begin{proof}
The map $\rho:S\to X$ is semismall, with the single relevant exceptional
fibre $E_{\exc}$.  Its top intersection form is the nondegenerate matrix
$(-2)$.  The semismall decomposition theorem
\cite[Theorem~3.4.1]{deCataldoMigliorini2002} gives
\begin{equation}
 R\rho_*\QQ_S[2]\simeq\IC_X\oplus\delta_0.
\label{eq:semismall}
\end{equation}
Taking Euler characteristics and using \cref{lem:surface-invariants} gives
\begin{equation}
 \chi(\IC_X)=e(S)-1=7.
\label{eq:chi-seven}
\end{equation}
A translation preserving $X$ must preserve its unique singular point, hence
\begin{equation}
 \Stab_P(X)=\{0\}.
\label{eq:stabilizer}
\end{equation}
In particular, the origin is the unique centre of the symmetry
\eqref{eq:symmetry}.
\end{proof}

Let $\Lambda_X\subset T^*P$ be the closure of the conormal bundle of
$X_{\reg}$, and let $\Lambda_0=T_0^*P$ be the conormal fibre at the origin.
Let
\[
 \gamma_X:\PP(\Lambda_X)\longrightarrow\PP(T_0^*P)
\]
be the map induced by cotangent projection.

\begin{lemma}[Conormal degree]\label[lemma]{lem:conormal-degree}
The characteristic cycle and conormal Gauss degree are
\[
 \operatorname{CC}(\IC_X)=\Lambda_X+\Lambda_0,
 \qquad
 \deg(\gamma_X)=6.
\]
In particular, the conormal Gauss map is dominant.
\end{lemma}

